\section{Funktionale Anforderungen}
Die funktionalen Anforderungen beschreiben alle Funktionen, die das System bereitstellen muss.
Diese ergeben sich direkt aus der Zielsetzung sowie aus den identifizierten Defiziten bestehender Systeme.

Die zentrale Funktion des Systems besteht in der kontinuierlichen und unterbrechungsfreien Darstellung der Uhrzeit in sprachlicher Form.
Im Gegensatz zu klassischen Wortuhren erfolgt zusätzlich eine minutengenaue und sekundengenaue Ergänzung, ohne die Hauptanzeige zu überlagern.
Weiterhin werden externe Datenquellen integriert.
Hierzu zählen insbesondere Wetterinformationen sowie die Außentemperatur.
Diese Daten werden in abstrahierter Form dargestellt, um die sprachliche Konsistenz der Anzeige zu erhalten.

Zusätzliche Sensorik ermöglicht die Erfassung von Innenraumparametern.
Die Anzeige erfolgt über separate \gls{led}-Elemente und ergänzt die Hauptfunktion um informationsorientierte Inhalte.
Die Benutzerinteraktion erfolgt über eine webbasierte Schnittstelle.
Dadurch wird eine flexible Steuerung ohne zusätzliche Hardware ermöglicht.
Alle funktionalen Anforderungen sind in Tabelle~\ref{tab:Funktionale Anforderung} aufgelistet.

\begin{table}[H]
\centering
\begin{tabular}{p{3cm} p{9cm}}
\textbf{ID} & \textbf{Beschreibung} \\
\hline
F1 & Wortbasierte Anzeige der aktuellen Uhrzeit (Stunde und Minute) \\
F2 & Kontinuierliche Sekundenanzeige ohne Unterbrechung von F1 \\
F3 & Minutengenaue Darstellung durch zusätzliche \gls{led}-Elemente \\
F4 & Anzeige der Außentemperatur in sprachlicher Form basierend auf Schwellenwerten \\
F5 & Darstellung von Wetterzuständen in sprachlicher Form \\
F6 & Anzeige der Innenraumtemperatur \\
F7 & Anzeige der Luftfeuchtigkeit im Innenraum \\
F8 & Symbolische Darstellung der Tageszeit \\
F9 & Symbolische Darstellung der Mondphase \\
F10 & Abruf externer Daten über Internetverbindung \\
F11 & Benutzerinteraktion zur Steuerung von Anzeige, Helligkeit und Farbe \\
F12 & Bereitstellung einer Oberfläche zur Systemsteuerung und Darstellung des Betriebszustand \\
\end{tabular}
\caption{Funktionale Anforderungen der Wortuhr}
\label{tab:Funktionale Anforderung}
\end{table}

%%%%%%%%%%%%%%%%%%%Old Version
% Welche Funktionen das System erfüllen muss
% Möglichst strukturiert (Tabelle)
% „Das System soll …“
% Die funktionalen Anforderungen beschreiben die vom System bereitzustellenden Funktionen. Sie definieren das erwartete Verhalten aus Sicht der Nutzer:innen.

% \begin{table}[H]
% \centering
% \begin{tabular}{|p{2cm}|p{11cm}|}
% \hline
% \textbf{ID} & \textbf{Anforderung} \\ \hline
% F1 & Das System soll die aktuelle Uhrzeit in sprachlicher Form als Wortanzeige darstellen. \\ \hline
% F2 & Das System soll die Uhrzeit mit einer Genauigkeit von mindestens 1 Minute anzeigen. \\ \hline
% F3 & Das System soll optional eine minutengenaue oder sekundengenaue Zusatzanzeige bereitstellen. \\ \hline
% F4 & Das System soll standortbezogene Wetterdaten anzeigen. \\ \hline
% F5 & Das System soll die aktuelle Mondphase anzeigen. \\ \hline
% F6 & Das System soll die Anzeigehelligkeit automatisch an die Umgebungshelligkeit anpassen. \\ \hline
% F7 & Das System soll eine manuelle Einstellung von Uhrzeit, Standort und Anzeigeparametern ermöglichen. \\ \hline
% F8 & Das System soll eine drahtlose Konfiguration über ein lokales Netzwerk unterstützen. \\ \hline
% F9 & Das System soll bei Stromausfall die Uhrzeit über eine geeignete Methode sichern. \\ \hline
% F10 & Das System soll verschiedene Anzeigezustände zyklisch oder nutzer:innenabhängig umschalten können. \\ \hline
% \end{tabular}
% \caption{Funktionale Anforderungen}
% \end{table}

% Die Anforderungen F1 und F2 definieren die Kernfunktion der Wortuhr. Die Anforderungen F3 bis F5 erweitern das System um Zusatzinformationen. Die Anforderungen F6 bis F10 erhöhen den Bedienkomfort sowie die Praxistauglichkeit im Alltag.